\documentclass[12pt]{article}
\usepackage[letterpaper,left=2.25cm,top=3.0cm,right=2.25cm,bottom=2.0cm, headsep=1.5cm]{geometry}
\usepackage{fancyhdr}
\usepackage[utf8]{inputenc}
\usepackage[T1]{fontenc}
\usepackage[spanish,mexico]{babel}
%\usepackage[letterpaper,margin=2.5cm]{geometry}
\usepackage{setspace}
\usepackage{titlesec}
\usepackage{tocloft}
\usepackage{graphicx}
\usepackage{array}

% Formato del documento
\usepackage[a4paper,margin=2.5cm]{geometry}
\usepackage{parskip}

% Tablas
\usepackage{float}
\usepackage{tabularx}
\usepackage{booktabs}
\usepackage{longtable}
\usepackage{enumitem}

% Imágenes

% Referencias y enlaces
\usepackage{xurl}
\usepackage[hidelinks]{hyperref}
\usepackage{csquotes}

% Referencias bibliográficas ISO 690
\usepackage{csquotes}

\usepackage[
    backend=biber,
    style=iso-numeric,
    sorting=none
]{biblatex}

\addbibresource{Referencias.bib}

% Interlineado
\onehalfspacing

% Ajustes de tablas
\renewcommand{\arraystretch}{1.3}

\usepackage[spanish,es-tabla]{babel}

\usepackage{csquotes}
\newcolumntype{C}[1]{>{\centering\arraybackslash}p{#1}} % Centrar 

%\usepackage{hyperref}
\usepackage{wallpaper}
% Para utilizar el formato de citas IEEE y comentar los dos parrafos siguientes. Ademas de utilizar el comando \cite en vez de \cite
\usepackage[backend=biber,style=iso-numeric]{biblatex}

% Estas lineas permiten romper los hipervinculos muy largos en las referencias!!!!
\setcounter{biburllcpenalty}{7000}
\setcounter{biburlucpenalty}{8000}
\addbibresource{Referencias.bib} % ARCHIVO DE BIBLIOGRAFÍA


\usepackage[bookmarks=true,breaklinks=true,bookmarksopen=false,colorlinks=true,linkcolor=blue]{hyperref}
\usepackage[hyphenbreaks]{breakurl}
\usepackage{pdfpages} % Incluir PDF en documento en LATEX
\setstretch{1.3}

\usepackage{soul}


% CONSTANTES NECESARIAS PARA EL DOCUMENTO ---> MODIFIQUEN A SU CRITERIO
\newcommand{\NombreGina}{Georgina Reta Limas}
\newcommand{\NombreJazmin}{Melissa Jazmin Torres Martinez}
\soulregister{\NombreGina}{1}
\soulregister{\NombreJazmin}{1}
\newcommand{\NombreTarea}{Prevención de extracción de fotografias de estudiantiles en entornos no cifrados (SIISU)} 
\newcommand{\NombreProyectoheader}{\NombreTarea}
\newcommand{\TipoProyecto}{Reporte Unidad 2}
\newcommand{\Materia}{Ética y legislación en tecnologías de la Información}
\newcommand{\Maestro}   {Dr. Héctor Hugo Avilés Arriaga}
\newcommand{\MatriculaGina}{2430043}
\newcommand{\MatriculaJazmin}{2430017}
\newcommand{\ncuatrimestre}{Séptimo}


\newcommand{\fechaRealizacion}{09 de Septiembre de 2026}
\newcommand{\fechaLimite}{11 de Septiembre de 2026}

\newcommand{\ncarrera}{Ingeniería en Tecnologías de la Información e Innovación Digital}


\fancyhead{} % Clear all header fields
\fancyhead[C]{%
    \begin{tabular}{@{} m{0.2\textwidth} m{0.5\textwidth} m{0.2\textwidth} @{}}
        % Left image
        \includegraphics[height=1.20cm]{UTYP.png}
        &
        \centering
        \scriptsize{\NombreProyectoheader}
        &
        \includegraphics[height=1.5cm]{LogoUPV_2023.png}
    \end{tabular}
}


\addto\captionsspanish{% Replace "english" with the language you use
  \renewcommand{\contentsname}%
    {Contenido temático}%
}


\renewcommand\spanishtablename{Tabla}

\begin{document}

% PAGINA 1 - PORTADA
\setcounter{page}{0}
\pagenumbering{roman}
\thispagestyle{empty}

\begin{titlepage}
\centering
\vspace*{-1.5cm}

% Logotipos superiores
\begin{tabular}{cp{6cm}c}
\includegraphics[height=2.25cm]{UTYP.png} & 
& \includegraphics[height=2.25cm]{LogoUPV_2023.png}   \\
\end{tabular}

\vspace{0.4cm}

{\Large \textbf{UNIVERSIDAD POLITÉCNICA DE VICTORIA}}\\[0.3cm]
{\large \MakeUppercase{\Materia}}\\[2cm]
{\Large \textbf{\MakeUppercase{\NombreTarea}}}

\vfill

% --- RECUADRO PARA IMAGEN CENTRAL ---
% Cambia 'tu_imagen.png' por el archivo deseado o ajusta 'height' si es necesario
%\includegraphics[width=6.5cm, height=6.5cm, keepaspectratio]{figs/portada.jpg}}
% ------------------------------------

\vfill

{\small QUE PRESENTA}\\[0.15cm]
\textbf{\MakeUppercase{\NombreGina}}\\[0.15cm]
\textbf{\MakeUppercase{\NombreJazmin}}\\[0.4cm]

{\small MATRICULAS}\\[0.15cm]
\textbf{\small{\MatriculaGina}}\\[0.15cm]
\textbf{\small{\MatriculaJazmin}}\\[0.4cm]

{\small EN CUMPLIMIENTO DEL DEBER}\\[0.15cm]
\textbf{\MakeUppercase{\TipoProyecto}}\\[0.4cm]

{\small MAESTRO A CARGO DE LA MATERIA}\\[0.15cm]
\textbf{\MakeUppercase{\Maestro}}

\vfill

\begin{flushright}
Ciudad Victoria, Tamaulipas, \fechaRealizacion.
\end{flushright}

\end{titlepage}

\tableofcontents
\clearpage

\pagenumbering{arabic}
\setcounter{page}{1}

\section{Introducción}

La digitalización de los procesos académicos ha permitido
a las instituciones de educación superior optimizar la
administración de información estudiantil mediante
plataformas digitales. Sin embargo, el almacenamiento
y tratamiento de datos personales exige implementar
mecanismos de seguridad que garanticen su confidencialidad,
integridad y disponibilidad.

La Universidad Politécnica de Victoria (UPV) utiliza
el Sistema Integral de Información y Servicios
Universitarios (SIISU) como una de sus principales
herramientas para la gestión académica. Esta plataforma
permite consultar horarios, calificaciones, historial
académico y datos personales, además de facilitar
actividades administrativas y docentes.

Durante la revisión del funcionamiento de SIISU,
se identificó que las fotografías de perfil de los
estudiantes, originalmente capturadas por la institución
para la elaboración de credenciales universitarias,
se encuentran asociadas a direcciones web que
contienen identificadores académicos predecibles.

De acuerdo con las observaciones realizadas, estos
recursos pueden visualizarse sin necesidad de
autenticación, lo que representa un riesgo potencial
de acceso no autorizado, recopilación de fotografías
y asociación de información personal.

Esta problemática adquiere relevancia ante los
antecedentes de utilización indebida de fotografías
estudiantiles en plataformas digitales, incluyendo
casos de exposición no consentida y manipulación
de imágenes mediante inteligencia artificial.

En consecuencia, el presente proyecto tiene como
objetivo analizar los riesgos jurídicos, técnicos y
éticos asociados con la exposición de fotografías
estudiantiles en SIISU, así como desarrollar una
política institucional de uso responsable de datos
personales y recursos digitales.

La propuesta contempla medidas preventivas,
controles de acceso, responsabilidades institucionales
y procedimientos de actuación orientados a
salvaguardar la privacidad de la comunidad
universitaria.

% ---------------------- HECHOS +++++++++++++++++++++++++++++
\section{Descripción de la organización y del contexto}
\subsection{Descripción de la institución}

La Universidad Politécnica de Victoria (UPV) es una
institución pública de educación superior ubicada
en Ciudad Victoria, Tamaulipas, México.

Como parte de sus actividades académicas y
administrativas, la universidad recopila,
almacena y procesa información personal de
estudiantes, docentes y personal administrativo.

Para gestionar estos procesos, dispone de
plataformas institucionales que permiten
administrar registros académicos, consultar
información estudiantil y facilitar la
comunicación entre los integrantes de la
comunidad universitaria.

Entre estas herramientas destaca SIISU,
una plataforma desarrollada por la propia
institución que constituye uno de los
principales medios de gestión y consulta
de información académica.

\subsection{Funcionamiento de SIISU}

SIISU es una plataforma institucional utilizada
por estudiantes y docentes para desarrollar
diversas actividades académicas.

En el caso de los estudiantes, permite consultar
horarios, historial académico, rendimiento,
calificaciones e información relacionada con
sus tutores.

Asimismo, dispone de un apartado de perfil
que contiene datos personales como nombres,
apellidos, correo electrónico, número telefónico
e información de referencia, además de
opciones para la administración de la cuenta.

En la parte superior del perfil se encuentra
la fotografía institucional del estudiante,
capturada previamente por la universidad
para la elaboración de su credencial.

Esta fotografía es administrada por la
institución y no puede ser modificada
directamente por el estudiante.

Por otra parte, los docentes utilizan SIISU
para registrar calificaciones y realizar
actividades relacionadas con la gestión
académica.

Durante la revisión que dio origen al
presente proyecto, la exposición descrita
se identificó en fotografías de estudiantes.
No se encontraron evidencias equivalentes
en los perfiles docentes examinados.

\subsection{Identificación de la problemática}

Durante una revisión exploratoria de SIISU,
se observó que las fotografías estudiantiles
se encuentran almacenadas mediante una
estructura de direcciones web que incorpora
la matrícula del estudiante.

El patrón identificado presenta la siguiente
estructura general:

\begin{center}
\small
\texttt{
/storage/documentos/\\
\{matricula\}/Foto/\{matricula\}.JPG
}
\end{center}

En esta estructura, la matrícula funciona
como identificador del directorio y del
archivo correspondiente a la fotografía.

El principal problema identificado no radica
únicamente en la utilización de identificadores
predecibles, sino en que, según las observaciones
realizadas, las imágenes pueden consultarse
sin necesidad de mantener una sesión
autenticada en SIISU.

Esta condición podría permitir que personas
ajenas al sistema accedan a fotografías
institucionales sin contar con autorización
para consultar dicha información.

Adicionalmente, la utilización de patrones
predecibles incrementa el riesgo de
recopilación automatizada de recursos,
particularmente cuando no existen
mecanismos suficientes de control de acceso,
supervisión y detección de solicitudes anómalas.

La situación adquiere mayor relevancia
considerando que las matrículas pueden
relacionarse con información disponible
en otros servicios institucionales,
facilitando la asociación entre
identificadores académicos e identidades
personales.

Por lo anterior, se considera necesario
evaluar los controles de protección de
fotografías estudiantiles y establecer
medidas que impidan su consulta,
recopilación o utilización no autorizada.

\begin{figure}[H]
    \centering
    \includegraphics[width=0.85\textwidth]
    {imagenes/esquema_siisu.jpg}
    \caption{Estructura general de acceso a
    fotografías estudiantiles en SIISU.}
    \label{fig:estructura_siisu}
\end{figure}

\subsection{Comparación con otros sistemas institucionales}

La Universidad Politécnica de Victoria
también dispone de SIITA, una plataforma
orientada principalmente a actividades
relacionadas con las tutorías académicas.

Durante la revisión exploratoria se
observaron diferencias entre los mecanismos
de identificación de fotografías utilizados
por SIISU y SIITA.

Mientras que SIISU presenta rutas que
incorporan directamente la matrícula,
SIITA aparentemente utiliza identificadores
que no permiten reconocer de forma
inmediata al estudiante mediante el
nombre del archivo.

Esta diferencia sugiere que la institución
ha implementado distintos mecanismos
de almacenamiento e identificación de
recursos digitales.

Sin embargo, la utilización de nombres
no reconocibles no garantiza por sí misma
que los archivos se encuentren cifrados
o protegidos mediante controles de acceso.

Por ello, se propone verificar los
mecanismos de autorización utilizados
en ambos sistemas e identificar
prácticas de seguridad que puedan
aplicarse de manera uniforme.

\subsection{Delimitación del problema}

El presente proyecto se enfoca en la
exposición de fotografías correspondientes
a perfiles estudiantiles de SIISU.

El análisis considera la estructura de
identificación de archivos, las condiciones
de acceso observadas y los riesgos derivados
de la posible consulta no autorizada.

No se contempla la extracción masiva de
fotografías, la enumeración automatizada
de matrículas ni la ejecución de pruebas
intrusivas sobre la infraestructura
universitaria.

Asimismo, no se afirma que los datos
académicos, las contraseñas o los perfiles
docentes presenten la misma condición
de exposición.

El alcance del proyecto consiste en
proponer una política institucional
de protección de datos personales,
acompañada de medidas técnicas,
jurídicas y administrativas que
permitan reducir los riesgos
identificados.



%-------------------------------------------------------------------------

\section{Identificación de datos personales
y recursos digitales involucrados}

\subsection{Identificación de datos personales}

La Universidad Politécnica de Victoria
administra información personal de sus
estudiantes mediante diferentes
plataformas institucionales.

Durante la revisión exploratoria se
identificaron dos niveles de exposición
de información: uno externo, relacionado
con las fotografías almacenadas en SIISU,
y otro interno, asociado con la
visibilidad de información estudiantil
en Google Workspace.

En el primer nivel, las fotografías
institucionales pueden consultarse,
según las observaciones realizadas,
sin necesidad de iniciar sesión.
Además, las direcciones utilizadas
incorporan directamente la matrícula
del estudiante.

En el segundo nivel, los usuarios
con cuentas institucionales pueden
localizar información de otros
integrantes mediante las herramientas
de búsqueda de Google Workspace.

Estas consultas pueden mostrar nombres
completos, direcciones de correo
institucional y matrículas. Asimismo,
se observó que determinados registros
del directorio permiten visualizar
números telefónicos.

Aunque cada dato puede cumplir
una finalidad administrativa legítima,
su disponibilidad conjunta incrementa
el riesgo de identificación, asociación
y utilización no autorizada de
información personal.

La Tabla \ref{tab:datos_personales}
presenta los principales datos
identificados y sus condiciones
de exposición.

\begin{table}[H]
\centering
\caption{Datos personales y condiciones de exposición}
\label{tab:datos_personales}
\small
\begin{tabularx}{\textwidth}{|X|X|X|}
\hline
\textbf{Dato personal} &
\textbf{Origen} &
\textbf{Accesibilidad observada} \\ \hline

Fotografía institucional &
SIISU &
Externa, sin autenticación. \\ \hline

Matrícula &
SIISU y correo institucional &
Visible en rutas y registros
institucionales. \\ \hline

Nombre completo &
Google Workspace &
Localizable por usuarios
institucionales. \\ \hline

Correo institucional &
Gmail y directorio &
Localizable por usuarios
institucionales. \\ \hline

Número telefónico &
Contactos institucionales &
Visible en determinados
registros internos. \\ \hline

Contraseña inicial &
Sistema de autenticación &
No se observó exposición
directa; existe un posible
riesgo por su estructura
predecible. \\ \hline

\end{tabularx}
\end{table}

\subsection{Recursos digitales involucrados}

Los recursos digitales involucrados
comprenden los sistemas institucionales
utilizados para gestionar, almacenar
y consultar información estudiantil.

El principal sistema analizado es SIISU,
debido a las condiciones observadas
en el acceso a fotografías.

Adicionalmente, Google Workspace
constituye un recurso relevante
por la posibilidad de consultar
información de los integrantes
de la comunidad universitaria.

\begin{table}[H]
\centering
\caption{Recursos digitales involucrados}
\label{tab:recursos_digitales}
\small
\begin{tabularx}{\textwidth}{|X|X|X|}
\hline
\textbf{Recurso} &
\textbf{Función institucional} &
\textbf{Riesgo identificado} \\ \hline

SIISU &
Gestión de información académica. &
Acceso público a fotografías. \\ \hline

Servidor de imágenes &
Almacenamiento de fotografías. &
Controles de acceso insuficientes. \\ \hline

Google Workspace &
Colaboración y administración
de usuarios. &
Correlación de datos personales. \\ \hline

Gmail institucional &
Comunicación académica. &
Visibilidad de identidades
y correos. \\ \hline

Google Contacts &
Directorio institucional. &
Exposición interna de
datos de contacto. \\ \hline

Sistema de autenticación &
Administración de credenciales. &
Posible predictibilidad de
contraseñas iniciales. \\ \hline

\end{tabularx}
\end{table}

\subsection{Riesgos derivados de la
correlación de información}

La exposición de datos personales
en diferentes plataformas puede
generar riesgos adicionales cuando
la información se relaciona
para construir perfiles individuales.

En el caso analizado, SIISU
presenta una posible exposición
externa de fotografías asociadas
con matrículas estudiantiles.

Por otra parte, las herramientas
de búsqueda institucional de
Google Workspace permiten
localizar información de contacto
e identificación de los usuarios. La combinación de ambas condiciones
podría facilitar la asociación
entre fotografías, matrículas,
nombres completos, correos
institucionales y números telefónicos.

Este fenómeno se conoce como
correlación de datos y representa
un riesgo para la privacidad,
especialmente cuando la información
puede utilizarse fuera de las
finalidades institucionales. Entre sus posibles consecuencias
se encuentran la suplantación
de identidad, el acoso,
la ingeniería social y
la utilización no consentida
de fotografías personales.

Es importante señalar que
la visibilidad de información
dentro de un directorio
institucional no constituye
automáticamente una vulnerabilidad.

Para determinar si existe
una exposición indebida,
es necesario evaluar la
finalidad del directorio,
los permisos establecidos
y la necesidad de que
cada dato sea accesible
para los distintos usuarios.

\subsection{Clasificación de los riesgos
según el nivel de acceso}

\begin{table}[H]
\centering
\caption{Clasificación de exposición
según el tipo de usuario}
\label{tab:niveles_exposicion}
\small
\begin{tabularx}{\textwidth}{|X|X|X|}
\hline
\textbf{Tipo de usuario} &
\textbf{Información potencialmente accesible} &
\textbf{Riesgo principal} \\ \hline

Usuario externo &
Fotografía y matrícula
asociada a la ruta. &
Acceso no autorizado
a imágenes. \\ \hline

Usuario institucional &
Fotografía, matrícula,
nombre, correo y
determinados teléfonos. &
Correlación de información
personal. \\ \hline

\end{tabularx}
\end{table}

La clasificación anterior
corresponde a los escenarios
de riesgo identificados.
No implica que se haya
comprobado el compromiso
de cuentas institucionales.


\section{Identificación de datos personales
y recursos digitales involucrados}

\subsection{Identificación de datos personales}

La Universidad Politécnica de Victoria administra
información personal de sus estudiantes mediante
diferentes plataformas institucionales.

De acuerdo con la Ley General de Protección de
Datos Personales en Posesión de Sujetos Obligados,
las instituciones públicas deben garantizar
el tratamiento adecuado y la protección de
la información personal bajo su responsabilidad
\parencite{lgpdppso2025}.

Durante la revisión exploratoria se identificaron
dos niveles de exposición de información:
uno externo, relacionado con las fotografías
almacenadas en SIISU, y otro interno,
asociado con la visibilidad de información
estudiantil en Google Workspace.

En el primer nivel, las fotografías institucionales
pueden consultarse, según las observaciones
realizadas, sin necesidad de iniciar sesión.
Además, las direcciones utilizadas incorporan
directamente la matrícula del estudiante.

Esta condición representa un posible problema
de control de acceso. OWASP establece que
las vulnerabilidades de referencia directa
insegura a objetos (IDOR) ocurren cuando
una aplicación permite acceder a recursos
sin verificar adecuadamente los permisos
del solicitante \parencite{owaspidor}.

En el segundo nivel, los usuarios con cuentas
institucionales pueden localizar información
de otros integrantes mediante las herramientas
de búsqueda de Google Workspace.

Estas consultas pueden mostrar nombres
completos, direcciones de correo institucional
y matrículas. Asimismo, se observó que
determinados registros del directorio
permiten visualizar números telefónicos.

Aunque cada dato puede cumplir una finalidad
administrativa legítima, su disponibilidad
conjunta incrementa el riesgo de identificación,
asociación y utilización no autorizada
de información personal.

La Tabla \ref{tab:datos_personales}
presenta los principales datos identificados
y sus condiciones de exposición.

\begin{table}[H]
\centering
\caption{Datos personales y condiciones de exposición}
\label{tab:datos_personales}
\small
\begin{tabularx}{\textwidth}{|X|X|X|}
\hline
\textbf{Dato personal} &
\textbf{Origen} &
\textbf{Accesibilidad observada} \\ \hline

Fotografía institucional &
SIISU &
Externa, sin autenticación. \\ \hline

Matrícula &
SIISU y correo institucional &
Visible en rutas y registros institucionales. \\ \hline

Nombre completo &
Google Workspace &
Localizable por usuarios institucionales. \\ \hline

Correo institucional &
Gmail y directorio &
Localizable por usuarios institucionales. \\ \hline

Número telefónico &
Contactos institucionales &
Visible en determinados registros internos. \\ \hline

Contraseña inicial &
Sistema de autenticación &
No se observó exposición directa;
existe un posible riesgo por su
estructura predecible. \\ \hline

\end{tabularx}
\end{table}


\subsection{Recursos digitales involucrados}

Los recursos digitales involucrados comprenden
los sistemas institucionales utilizados para
gestionar, almacenar y consultar información
estudiantil.

El principal sistema analizado es SIISU,
debido a las condiciones observadas en
el acceso a fotografías.

Adicionalmente, Google Workspace constituye
un recurso relevante por la posibilidad
de consultar información de los integrantes
de la comunidad universitaria.

Google Workspace dispone de herramientas
de directorio que permiten localizar
información de usuarios dentro de una
organización. La visibilidad de estos
datos depende de la configuración
administrativa y de los permisos
establecidos \parencite{googleworkspace}.

\begin{table}[H]
\centering
\caption{Recursos digitales involucrados}
\label{tab:recursos_digitales}
\small
\begin{tabularx}{\textwidth}{|X|X|X|}
\hline
\textbf{Recurso} &
\textbf{Función institucional} &
\textbf{Riesgo identificado} \\ \hline

SIISU &
Gestión de información académica. &
Acceso público a fotografías. \\ \hline

Servidor de imágenes &
Almacenamiento de fotografías. &
Controles de acceso insuficientes. \\ \hline

Google Workspace &
Colaboración y administración de usuarios. &
Correlación de datos personales. \\ \hline

Gmail institucional &
Comunicación académica. &
Visibilidad de identidades y correos. \\ \hline

Google Contacts &
Directorio institucional. &
Exposición interna de datos de contacto. \\ \hline

Sistema de autenticación &
Administración de credenciales. &
Posible predictibilidad de
contraseñas iniciales. \\ \hline

\end{tabularx}
\end{table}


\subsection{Riesgos derivados de la
correlación de información}

La exposición de datos personales en
diferentes plataformas puede generar
riesgos adicionales cuando la información
se relaciona para construir perfiles
individuales.

En el caso analizado, SIISU presenta
una posible exposición externa de
fotografías asociadas con matrículas
estudiantiles.

Por otra parte, las herramientas de
búsqueda institucional de Google Workspace
permiten localizar información de
contacto e identificación de los usuarios.

La combinación de ambas condiciones
podría facilitar la asociación entre
fotografías, matrículas, nombres completos,
correos institucionales y números telefónicos.

Este fenómeno se relaciona con los riesgos
de privacidad derivados de la agregación
y vinculación de información personal.

El NIST Privacy Framework reconoce
que el procesamiento y la asociación
de datos pueden generar riesgos para
las personas, por lo que recomienda
identificarlos y gestionarlos mediante
controles organizacionales y técnicos
\parencite{nistprivacy}.

Entre las posibles consecuencias de
esta correlación se encuentran la
suplantación de identidad, el acoso,
la ingeniería social y la utilización
no consentida de fotografías personales.

Es importante señalar que la visibilidad
de información dentro de un directorio
institucional no constituye automáticamente
una vulnerabilidad.

Para determinar si existe una exposición
indebida, es necesario evaluar la finalidad
del directorio, los permisos establecidos
y la necesidad de que cada dato sea
accesible para los distintos usuarios.


\subsection{Riesgos relacionados con
las credenciales de acceso}

Durante la revisión del contexto
institucional se identificó una posible
debilidad relacionada con la generación
de contraseñas iniciales para las
cuentas estudiantiles.

De acuerdo con la información recopilada,
estas credenciales siguen una estructura
parcialmente predecible basada en
información personal del estudiante.

Asimismo, se señaló que algunos usuarios
conservan la contraseña asignada
originalmente por la institución.

Esta práctica podría incrementar el riesgo
de acceso no autorizado cuando las
credenciales iniciales no son suficientemente
aleatorias o permanecen activas durante
periodos prolongados.

El National Institute of Standards and
Technology (NIST) recomienda utilizar
mecanismos de autenticación resistentes
a ataques de adivinación y evitar
prácticas que debiliten la seguridad
de las credenciales
\parencite{nist80063b}.

Sin embargo, no se cuenta con evidencia
de que las contraseñas se encuentren
públicamente expuestas ni de que se
hayan comprometido cuentas mediante
este mecanismo.

Como medida preventiva, se recomienda
implementar contraseñas iniciales
aleatorias, de uso temporal, y exigir
su modificación durante el primer
inicio de sesión.

Adicionalmente, se propone incorporar
mecanismos de autenticación multifactor
y procedimientos seguros de recuperación
de cuentas.


\subsection{Clasificación de los riesgos
según el nivel de acceso}

La exposición de información presenta
diferencias según el tipo de usuario
y los permisos disponibles.

Esta clasificación permite distinguir
los riesgos asociados con el acceso
externo a fotografías de aquellos
relacionados con la visibilidad
de información dentro del entorno
institucional.

\begin{table}[H]
\centering
\caption{Clasificación de exposición
según el tipo de usuario}
\label{tab:niveles_exposicion}
\small
\begin{tabularx}{\textwidth}{|X|X|X|}
\hline
\textbf{Tipo de usuario} &
\textbf{Información potencialmente accesible} &
\textbf{Riesgo principal} \\ \hline

Usuario externo &
Fotografía y matrícula
asociada a la ruta. &
Acceso no autorizado a imágenes. \\ \hline

Usuario institucional &
Fotografía, matrícula, nombre,
correo y determinados teléfonos. &
Correlación de información personal. \\ \hline

Usuario con credenciales comprometidas &
Información disponible según
los permisos de la cuenta afectada. &
Acceso indebido y suplantación
de identidad. \\ \hline

\end{tabularx}
\end{table}

La clasificación anterior corresponde
a los escenarios de riesgo identificados
y no implica que se haya comprobado
el compromiso de cuentas institucionales.

De acuerdo con OWASP, los sistemas deben
aplicar controles de autorización para
cada solicitud y seguir el principio
del menor privilegio, otorgando únicamente
los permisos necesarios para cada usuario
\parencite{owaspauth}.

\section{Política de uso responsable de datos
personales y recursos digitales}

\subsection{Propósito}

Establecer lineamientos institucionales para
garantizar la protección de fotografías
estudiantiles y demás datos personales
almacenados en SIISU, mediante mecanismos
de seguridad, control de acceso y
tratamiento responsable de la información.

\subsection{Alcance}

La presente política está dirigida
al personal administrativo, docentes,
desarrolladores, responsables de
tecnologías de la información y
demás integrantes de la comunidad
universitaria que intervengan en
el tratamiento de datos personales.

Su aplicación comprende los sistemas
institucionales, servidores, bases de
datos y servicios digitales utilizados
para almacenar o consultar información
estudiantil.

\subsection{Tratamiento de datos personales}

La institución deberá garantizar
que los datos personales sean
tratados conforme a las finalidades
para las cuales fueron recopilados.

Se establecen los siguientes lineamientos:

\begin{enumerate}
    \item Recopilar únicamente la
    información necesaria para
    las actividades institucionales.

    \item Informar a los titulares
    sobre las finalidades del
    tratamiento mediante el
    aviso de privacidad.

    \item Restringir el acceso a
    fotografías y registros
    personales a usuarios autorizados.

    \item Prohibir la difusión
    de fotografías institucionales
    para fines no autorizados.

    \item Conservar la información
    durante los periodos legal
    y administrativamente aplicables.

    \item Eliminar o bloquear los
    datos cuando corresponda,
    conforme a la normativa vigente.
\end{enumerate}

\subsection{Medidas de privacidad y seguridad}

El área responsable de tecnologías
de la información deberá implementar
medidas preventivas para proteger
las fotografías institucionales.

Estas medidas comprenderán:

\begin{enumerate}
    \item Exigir autenticación para
    consultar fotografías estudiantiles.

    \item Verificar la autorización
    del usuario para acceder
    a cada recurso solicitado.

    \item Evitar la exposición directa
    de archivos personales mediante
    directorios públicos.

    \item Utilizar identificadores
    no predecibles como medida
    complementaria de seguridad.

    \item Implementar registros
    de acceso y mecanismos de
    detección de actividad anómala.

    \item Establecer restricciones
    frente a solicitudes automatizadas
    excesivas.

    \item Aplicar cifrado durante
    la transmisión de información
    y, cuando corresponda,
    durante su almacenamiento.

    \item Realizar evaluaciones
    periódicas de seguridad.
\end{enumerate}

Estas medidas se fundamentan
en las recomendaciones de OWASP
para prevenir vulnerabilidades
de referencia directa insegura
a objetos \parencite{owaspidor}.

\subsection{Derechos ARCO}

La universidad deberá disponer
de procedimientos para que
los titulares ejerzan sus
derechos de acceso,
rectificación, cancelación
y oposición, conforme
a las disposiciones aplicables.

Las solicitudes deberán
canalizarse a la Unidad
de Transparencia o instancia
institucional competente.

El procedimiento deberá
contemplar la recepción,
verificación de identidad,
evaluación de procedencia
y emisión de una respuesta
dentro de los plazos legales.

El ejercicio de estos derechos
estará sujeto a las condiciones,
excepciones y obligaciones
de conservación previstas
en la legislación.

\subsection{Uso de recursos digitales de terceros}

El personal institucional
deberá utilizar software,
imágenes, textos, videos,
bases de datos y demás
recursos digitales respetando
las condiciones de sus
respectivas licencias.

Antes de incorporar recursos
de terceros, deberá verificarse:

\begin{enumerate}
    \item La procedencia del recurso.

    \item La titularidad de los
    derechos correspondientes.

    \item La existencia de licencias
    o autorizaciones válidas.

    \item Los permisos de reproducción,
    modificación y distribución.

    \item Las obligaciones de
    atribución y reconocimiento.

    \item Las restricciones para
    usos comerciales o institucionales.
\end{enumerate}

La disponibilidad pública
de una fotografía no implica
que exista autorización para
su reutilización.

\subsection{Responsabilidades institucionales}

\begin{table}[H]
\centering
\caption{Responsabilidades propuestas}
\small
\begin{tabularx}{\textwidth}{|l|X|}
\hline
\textbf{Responsable} &
\textbf{Funciones} \\ \hline

Dirección institucional &
Aprobar y supervisar la política. \\ \hline

Área de TI &
Implementar controles de seguridad. \\ \hline

Desarrolladores &
Corregir fallas y verificar permisos. \\ \hline

Unidad de Transparencia &
Atender solicitudes ARCO. \\ \hline

Área jurídica &
Verificar obligaciones y licencias. \\ \hline

Docentes y administrativos &
Resguardar información estudiantil. \\ \hline

Estudiantes &
Utilizar responsablemente los sistemas
y reportar incidentes. \\ \hline

\end{tabularx}
\end{table}

\subsection{Actuación ante incidentes}

Cuando se identifique una posible
exposición de fotografías o
información personal, deberán
realizarse las siguientes acciones:

\begin{enumerate}
    \item Registrar y comunicar
    el incidente mediante los
    canales institucionales.

    \item Evaluar el alcance
    de la exposición.

    \item Restringir temporalmente
    los accesos inseguros.

    \item Preservar las evidencias
    técnicas necesarias sin
    difundir datos personales.

    \item Corregir las condiciones
    que originaron el incidente.

    \item Determinar las obligaciones
    de notificación a titulares
    y autoridades competentes.

    \item Verificar la efectividad
    de las medidas correctivas.

    \item Documentar las acciones
    realizadas y dar seguimiento.
\end{enumerate}

\subsection{Evidencias de cumplimiento}

Para comprobar la aplicación
de la política, deberán
conservarse evidencias como:

\begin{enumerate}
    \item Registros de acceso
    a recursos protegidos.

    \item Informes de auditorías
    de seguridad autorizadas.

    \item Documentación de controles
    de autenticación y autorización.

    \item Inventarios de datos
    personales y recursos digitales.

    \item Registros de atención
    de solicitudes ARCO.

    \item Reportes de incidentes
    y acciones correctivas.

    \item Inventarios de licencias
    y autorizaciones de terceros.

    \item Constancias de capacitación
    del personal.
\end{enumerate}

\section{Decálogo de actuación institucional}

\begin{enumerate}[label=\textbf{\arabic*.}]

    \item \textbf{Proteger la identidad estudiantil.}
    Las fotografías y datos personales
    deberán tratarse como información
    de acceso restringido.

    \item \textbf{Verificar cada acceso.}
    Ningún recurso personal podrá
    consultarse sin la autorización
    correspondiente.

    \item \textbf{Evitar identificadores expuestos.}
    Los sistemas deberán reducir
    la utilización de referencias
    predecibles a archivos personales.

    \item \textbf{Respetar la finalidad.}
    Las fotografías institucionales
    solo podrán utilizarse para
    los propósitos autorizados.

    \item \textbf{Prevenir la recopilación indebida.}
    Se deberán implementar controles
    para detectar accesos anómalos
    y solicitudes automatizadas.

    \item \textbf{Garantizar los derechos ARCO.}
    Los titulares podrán ejercer
    sus derechos conforme a
    los procedimientos legales.

    \item \textbf{Respetar los derechos de autor.}
    Los recursos digitales de
    terceros deberán utilizarse
    conforme a sus licencias.

    \item \textbf{Reportar incidentes.}
    Toda exposición o utilización
    indebida deberá comunicarse
    mediante los canales establecidos.

    \item \textbf{Asignar responsabilidades.}
    Cada área deberá conocer
    sus obligaciones relacionadas
    con la protección de información.

    \item \textbf{Evaluar y mejorar.}
    La seguridad de los sistemas
    deberá revisarse periódicamente
    para prevenir nuevas vulnerabilidades.

\end{enumerate}

\section{Justificación de las reglas}

La política propuesta responde a la
necesidad de garantizar un tratamiento
responsable de la información personal
administrada por la Universidad
Politécnica de Victoria.

Desde la perspectiva jurídica,
se fundamenta en el derecho
a la protección de datos personales,
reconocido por la Constitución
Política de los Estados Unidos
Mexicanos y desarrollado mediante
la legislación aplicable a
los sujetos obligados
\parencite{lgpdppso2025}.

Desde la perspectiva técnica,
las recomendaciones de OWASP
establecen que los sistemas
deben verificar los permisos
de acceso a los recursos,
independientemente de la
complejidad de sus identificadores
\parencite{owaspidor}.

Desde la perspectiva ética,
la institución tiene la
responsabilidad de prevenir
que información recopilada
con fines académicos sea
utilizada para propósitos
distintos a los autorizados.

\begin{table}[H]
\centering
\caption{Fundamentación de las reglas propuestas}
\label{tab:fundamentacion}
\small
\begin{tabularx}{\textwidth}{|c|X|X|}
\hline
\textbf{Regla} &
\textbf{Justificación} &
\textbf{Fundamento} \\ \hline

1 y 4 &
Protección de identidad y finalidad. &
Legislación de datos personales. \\ \hline

2 y 3 &
Prevención de accesos indebidos. &
OWASP IDOR. \\ \hline

5 &
Prevención de recopilación automatizada. &
Seguridad y control de acceso. \\ \hline

6 &
Garantía de derechos de los titulares. &
Legislación ARCO. \\ \hline

7 &
Uso legítimo de recursos de terceros. &
Derechos de autor y licencias. \\ \hline

8 &
Atención de incidentes de seguridad. &
Gestión de riesgos. \\ \hline

9 &
Asignación de obligaciones. &
Responsabilidad institucional. \\ \hline

10 &
Verificación y mejora continua. &
Buenas prácticas de seguridad. \\ \hline

\end{tabularx}
\end{table}

\subsection{Medidas correctivas prioritarias}

\begin{table}[H]
\centering
\caption{Plan de implementación propuesto}
\label{tab:implementacion}
\small
\begin{tabularx}{\textwidth}{|X|X|l|}
\hline
\textbf{Medida} &
\textbf{Responsable} &
\textbf{Prioridad} \\ \hline

Restringir el acceso público
a fotografías &
Área de TI &
Inmediata \\ \hline

Implementar autorización
por usuario y recurso &
Desarrolladores &
Alta \\ \hline

Revisar directorios públicos &
Área de TI &
Alta \\ \hline

Incorporar registros
y alertas de acceso &
Administradores &
Alta \\ \hline

Revisar mecanismos de SIITA &
Desarrolladores &
Media \\ \hline

Actualizar la política
de privacidad &
Área jurídica &
Alta \\ \hline

Realizar auditorías periódicas &
Área de TI &
Continua \\ \hline

Capacitar al personal &
Administración &
Continua \\ \hline

\end{tabularx}
\end{table}

\section{Conclusiones}

El análisis realizado permitió
identificar una posible deficiencia
de control de acceso en el
almacenamiento y distribución
de fotografías estudiantiles
mediante SIISU.

La utilización de direcciones
predecibles asociadas a matrículas,
combinada con la posibilidad
reportada de consultar imágenes
sin autenticación, representa
un riesgo para la privacidad
de la comunidad universitaria.

Los antecedentes documentados
demuestran que las fotografías
estudiantiles pueden ser
utilizadas para finalidades
distintas a aquellas para
las cuales fueron recopiladas.

Por ello, se propone una
política institucional que
establezca responsabilidades,
controles de autorización,
medidas preventivas y
procedimientos de actuación
ante incidentes.

La implementación de estas
disposiciones permitiría
fortalecer la protección
de datos personales,
reducir los riesgos
de exposición y promover
una cultura de seguridad
digital dentro de la
Universidad Politécnica
de Victoria.



\newpage

\printbibliography[
    heading=bibintoc,
    title={Referencias}
]
\end{document}